%%%%%%%%%%%%%%%%%%%%%%%%%%%%%%%%%
%
%       EXERCÍCIO
%
%%%%%%%%%%%%%%%%%%%%%%%%%%%%%%%%%

\definevalorquestao

\begin{exercicioBanco}[\valorquestao]
Analise as afirmações a respeito de potências e funções exponenciais:

\begin{enumerate}[label=\Roman*.]
    \item A função \(f(x) = \left(\sqrt{2}\right)^x\) é crescente em \(\mathbb{R}\).
    \item Se \(f(x) = e^x\), então \(f(a \cdot b) = f(a) \cdot f(b)\) para quaisquer \(a, b \in \mathbb{R}\).
    \item O valor numérico da expressão \(4^{-0{,}5} + 0{,}1^{-1}\) é igual a \(10{,}5\).
    \item Se \(m, n \in \mathbb{N} = \{1,2,3,\dots\}\) com \(m < n\), então \(a^m < a^n\) para todo \(a \in \mathbb{R}\).
\end{enumerate}

Assinale a alternativa correta:

% Define as alternativas
\newcommand{\alternativas}{%
    \begin{center}
        \begin{tabularx}{\textwidth}{XXX}
            (a) Todas são falsas. &
            (b) Todas são verdadeiras. &
            (c) Apenas I e III são verdadeiras. \\[5pt]
            (d) Apenas II e IV são verdadeiras. &
            (e) Apenas I, III e IV são verdadeiras.
        \end{tabularx}
    \end{center}
}

% Define a resposta correta
\newcommand{\resposta}{C}

% Lógica condicional para exibição
\ifdefstring{\atividade}{lista}{%
    \alternativas
    \vspace{0.5em}
    
    \noindent\textbf{Resposta:} letra \textbf{\resposta}.
}{%
    \ifdefstring{\modo}{objetivo}{%
        \alternativas
    }{%
        % modo = discursiva → não mostra alternativas
    }
}
\end{exercicioBanco}

